\begin{frame}{Metodologia}
  \begin{itemize}
    \item[$\blacktriangleright$] Conjuntos de dados
  \end{itemize}

  \vspace{-0.4em}
  \footnotesize
  \setlength{\tabcolsep}{3pt}
  \renewcommand{\arraystretch}{1.15}

  \begin{table}[h]
    \centering
    \begin{tabular}{|l|c|c|c|}
      \hline
      \textbf{Indicador} & \textbf{Lending Club} & \textbf{Taiwan} & \textbf{Corporate Credit} \\
      \hline
      Origem/Tipo & P2P (EUA) & Cartões (PF) & Empresas \\
      \hline
      Amostras & 630 mil & 30 mil & 2.029 \\
      \hline
      Atributos & 145 & 24 & 31 \\
      \hline
      Inadimplentes & 1,1\% & 22\% & 2,8\% \\
      \hline
    \end{tabular}
  \end{table}
  \normalsize
\end{frame}

\begin{frame}{Metodologia}
  \begin{itemize}
    \item[$\blacktriangleright$] Modelos de aprendizado de máquina
      \begin{itemize}
        \item SVM
        \item Random Forest
        \item XGBoost
        \item MLP
      \end{itemize}
  \end{itemize}
\end{frame}

\begin{frame}{Metodologia}
  \begin{itemize}
    \item[$\blacktriangleright$] Estratégias para lidar com desbalanceamento de dados
      \begin{itemize}
        \item SMOTE (oversampling)
        \item RUS (under-sampling)
        \item SMOTE-Tomek (híbrido)
        \item CS-SVM (sensível a custos)
      \end{itemize}
  \end{itemize}
\end{frame}

\begin{frame}{Metodologia}
  \begin{itemize}
    \item[$\blacktriangleright$] Métricas de desempenho
      \begin{itemize}
        \item Acurácia balanceada
        \item F1-score
        \item G-mean
        \item Precisão
        \item Sensibilidade
      \end{itemize}
  \end{itemize}
\end{frame}

\begin{frame}{Metodologia}
  \begin{itemize}
    \item[$\blacktriangleright$] Experimentos
      \begin{itemize}
        \item Validação cruzada estratificada
        \item Separação dos dados em conjuntos de treinamento e teste
          \begin{itemize}
            \item Treinamento: 70\%
            \item Teste: 30\%
          \end{itemize}
        \item Busca em grade para otimização de hiper-parâmetros
        \item Média de 5 sementes para cada combinação de modelo e estratégia, totalizando 85 experimentos por conjunto de dados.
      \end{itemize}
  \end{itemize}
\end{frame}
